\section{Введение}
Методы Франка-Вульфа заменяют проекцию линейной минимизацией на
компактном выпуклом множестве. Для выпуклой функции полный зазор
ограничивает глобальную ошибку. Для невыпуклой функции малый зазор
обычно означает лишь близость к стационарности. Поэтому для остановки
с гарантией глобальной точности нужна количественная связь наблюдений
с ошибкой.

Мы изучаем эту связь при сдвинутом RSI и границе
кривизны через значение ошибки:
\[
 \|\nabla^2 f(x)\|\le H_0+H_1\bigl(f(x)-f(x^*)\bigr).
\]
Минимум может лежать на границе и иметь ненулевой градиент.
Норма гессиана не возводится в квадрат. Такая модель используется
в анализе разогрева шага \cite{warmup}. Она отличается от границы
через $L_0+L_1\|\nabla f(x)\|$, изученной для метода Франка-Вульфа
в \cite{vyguzov}. Регулярность относительно фиксированного минимума
и условия секущей служат известными заменами выпуклости
\cite{regregularity,rsipep}.

Основной результат даёт точную границу ошибки по полному зазору $g$,
норме градиента $P$ и, если он известен, диаметру $D$.
Одно радиальное сравнение выводит отдельные оценки по норме и зазору,
а также совместные границы $\mathcal J(g,P)$ и $\mathcal J_D(g,P)$.
Они точны в классе функций и областей, удовлетворяющих условиям,
в том числе при точном задании обоих наблюдений и диаметра,
если $g\le PD$.
Точность для каждой заранее фиксированной области не утверждается.

Также найден точный критерий того, когда обычный зазор ограничивает
ошибку. Его проверка требует арифметики и квадратного корня.
В этом режиме совместная граница сводится к одному скалярному
уравнению; алгебраическая верхняя оценка обходится и без него.
Вне этого режима конечная сетка и бисекция дают гарантированные
включения с явным числом сравнений. Профиль по радиусу не обязан
быть вогнутым. Общее ограждение не предполагает ни вогнутости,
ни одновершинности.

Второй результат превращает подходящую нелинейную связь ошибки
и зазора в правило остановки с адаптивной точностью направленных
наблюдений. Счёт включает неудачные уточнения, отклонённые пробы
функции и последний запрос. Результат распространяется на шаги
от активных атомов и медианы средних с общей вероятностной гарантией для
всех адаптивных запросов. Итерация использует известный подбор
шага \cite{pedregosa,lacostejaggi}; здесь изучаются сертификаты
и полный счёт работы. Оракулы с относительной ошибкой и адаптивные
варианты метода Франка-Вульфа также исследованы в
\cite{oraclehallak,oracledenisov}.

В конце статьи приведены явные константы для реализуемой сигмоидной
регрессии на $\ell_1$-шаре и фиксированные проверки формул
сертификатов. Сами по себе доказанные оценки не устанавливают
ускорение по времени.
